\documentclass[9pt]{beamer}
\usetheme{metropolis}
\usepackage{amsmath}
\usepackage{amssymb}
\usepackage{graphicx}
\usepackage{booktabs}
\usepackage{xcolor}
\usepackage{listings}

% Define color palette matching our project design system
\definecolor{primary}{HTML}{1abc9c}    % Deep Teal
\definecolor{secondary}{HTML}{3498db}  % Astro Blue
\definecolor{darkbg}{HTML}{2c3e50}     % Slate Blue
\definecolor{accent}{HTML}{e67e22}     % Orange

\setbeamercolor{palette primary}{bg=darkbg, fg=white}
\setbeamercolor{palette secondary}{bg=primary, fg=white}
\setbeamercolor{palette tertiary}{bg=secondary, fg=white}
\setbeamercolor{structure}{fg=primary}
\setbeamercolor{alerted text}{fg=accent}

\title{Stellar Trace: Statistical Inference Foundation}
\subtitle{Meet 1 Stats Meet: From Probability Basics to 1D \& 2D Hypothesis Testing}
\author{Astrophysics Mentorship Program}
\date{\today}

\begin{document}

\maketitle

% Slide 2: Pedagogical Roadmap
\begin{frame}{Where We Are \& Where We Are Going}
  \begin{block}{Our Starting Point}
    So far, you have learned the basic axioms of probability, discrete events, conditional probability, and Bayes' Theorem:
    $$P(A \mid B) = \frac{P(B \mid A)P(A)}{P(B)}$$
  \end{block}
  \begin{block}{Today's Educational Objective}
    To bridge the gap between abstract probability theory and real-world astrophysics data. We want to answer:
    \textit{Do supernova host galaxies differ physically from standard field galaxies, or are they drawn from the same population?}
    To solve this, we will build a complete mathematical bridge:
    $$\text{Probability Recap} \rightarrow \text{PDFs} \rightarrow \text{CDFs} \rightarrow \text{eCDFs} \rightarrow \text{Hypothesis Testing} \rightarrow \text{1D KS/AD Tests} \rightarrow \text{2D KS Test}$$
  \end{block}
\end{frame}

% Slide 3: Meet Objectives
\begin{frame}{Meet Objectives}
  By the end of this meet, you will be able to:
  \begin{enumerate}
    \item Convert continuous Probability Density Functions (PDFs) to Cumulative Distribution Functions (CDFs) via integration.
    \item Construct Empirical Cumulative Distribution Functions (eCDFs) from raw telescope datasets.
    \item Formulate Null ($H_0$) and Alternative ($H_1$) hypotheses, test statistics, and significance levels ($\alpha$).
    \item Differentiate between Type I and Type II errors and understand their statistical trade-offs.
    \item Explain the mathematical mechanics, strengths, and tail-blindness limitations of the 1D Kolmogorov-Smirnov (K-S) test.
    \item Implement the Anderson-Darling (A-D) test using a tail-weighting function to analyze distribution extremes.
    \item Generalize hypothesis testing to 2D space using the quadrant-splitting Fasano \& Franceschini (1987) K-S test.
  \end{enumerate}
\end{frame}

% Slide 4: Section 1 Divider
\begin{frame}[plain]
  \vfill
  \centering
  \begin{beamercolorbox}[sep=8pt,center,shadow=true,rounded=true]{palette primary}
    \usebeamerfont{title}\textbf{Section 1: Probability Recap \& Continuous Variables}
  \end{beamercolorbox}
  \vfill
\end{frame}

% Slide 5a: Review of Basic Probability Rules - Foundations
\begin{frame}{Review of Basic Probability Rules: Foundations}
  Let's review the foundations of probability:
  \begin{itemize}
    \item \textbf{Sample Space ($S$):} The set of all possible outcomes (e.g., all galaxies in the universe).
    \item \textbf{Event ($A$):} A subset of the sample space (e.g., galaxies hosting a Core-Collapse Supernova).
    \item \textbf{Conditional Probability:} The probability of event $A$ occurring given that event $B$ has occurred:
      $$P(A \mid B) = \frac{P(A \cap B)}{P(B)} \quad (\text{for } P(B) > 0)$$
  \end{itemize}
\end{frame}

% Slide 5b: Review of Basic Probability Rules - The Product Rule
\begin{frame}{Review of Basic Probability Rules: The Product Rule}
  \begin{itemize}
    \item \textbf{The Product Rule:} Re-arranging conditional probability gives the joint probability of both events:
      $$P(A \cap B) = P(A \mid B)P(B) = P(B \mid A)P(A)$$
  \end{itemize}
  \begin{exampleblock}{Astronomical Intuition}
    Let $A$ be the event ``galaxy has high star formation rate'' and $B$ be the event ``galaxy is a spiral''. $P(A \mid B)$ is the probability that a spiral galaxy has a high star formation rate, which is higher than the average galaxy probability $P(A)$.
  \end{exampleblock}
\end{frame}

% Slide 6a: Bayes' Theorem in Action - Formulation
\begin{frame}{Bayes' Theorem in Action: Formulation}
  \begin{block}{Mathematical Formulation}
    Derived directly from the product rule:
    $$P(A \mid B) = \frac{P(B \mid A)P(A)}{P(B)}$$
    Where:
    \begin{itemize}
      \item $P(A \mid B)$ is the \textbf{Posterior Probability} (updated belief).
      \item $P(B \mid A)$ is the \textbf{Likelihood} (probability of observing $B$ if $A$ is true).
      \item $P(A)$ is the \textbf{Prior Probability} (initial belief).
      \item $P(B)$ is the \textbf{Evidence} (total probability of observing $B$):
        $$P(B) = P(B \mid A)P(A) + P(B \mid \neg A)P(\neg A)$$
    \end{itemize}
  \end{block}
\end{frame}

% Slide 6b: Bayes' Theorem in Action - Application
\begin{frame}{Bayes' Theorem in Action: Application}
  \begin{exampleblock}{Supernova Example}
    We observe a supernova in a red, passive galaxy. What is the probability it is Type Ia ($A$) vs. Core-Collapse ($\neg A$)?
    CC-SNe require massive star progenitors which are absent in passive galaxies, so $P(\text{Red} \mid \text{CC}) \approx 0$.
    Bayes' Theorem mathematically shows:
    $$P(\text{Ia} \mid \text{Red}) = \frac{P(\text{Red} \mid \text{Ia})P(\text{Ia})}{P(\text{Red} \mid \text{Ia})P(\text{Ia}) + P(\text{Red} \mid \text{CC})P(\text{CC})} \approx 1.0$$
  \end{exampleblock}
\end{frame}

% Slide 7: Continuous Random Variables
\begin{frame}{Continuous Random Variables}
  Discrete variables take countable values (e.g., number of supernovae = 0, 1, 2...). Continuous random variables can take any value within a continuous real interval.
  \begin{itemize}
    \item \textbf{Astrophysical Examples:}
      \begin{itemize}
        \item Stellar Mass: $\log_{10} M_\star \in [6.0, 12.0] \, M_\odot$
        \item Specific Star Formation Rate: $\log_{10} \text{sSFR} \in [-14.0, -8.0] \, \text{yr}^{-1}$
        \item Redshift: $z \in [0.0, 10.0]$
      \end{itemize}
    \item \textbf{Transition from Discrete to Continuous:}
      \begin{itemize}
        \item For discrete variables, we use a Probability Mass Function (PMF), where $P(X = x) = p(x)$.
        \item For continuous variables, as the spacing between possible values approaches zero, the PMF transitions to a continuous \textbf{Probability Density Function (PDF)}, $f(x)$.
      \end{itemize}
  \end{itemize}
\end{frame}

% Slide 8: Mathematical Axioms of PDFs
\begin{frame}{Mathematical Axioms of PDFs}
  A function $f(x)$ is a valid Probability Density Function (PDF) if and only if it satisfies:
  \begin{enumerate}
    \item \textbf{Non-negativity:} $f(x) \ge 0$ for all $x \in \mathbb{R}$. (Negative probability is physically meaningless).
    \item \textbf{Normalization:} The total area under the curve must be exactly 1:
      $$\int_{-\infty}^{\infty} f(x) \, dx = 1$$
    \item \textbf{Probability of an Interval:} The probability that $X$ falls between $a$ and $b$ is the area under $f(x)$ between those bounds:
      $$P(a \le X \le b) = \int_{a}^{b} f(x) \, dx$$
  \end{enumerate}
  \begin{alertblock}{The Zero-Probability Paradox}
    For a continuous variable, the probability of taking an exact single value is zero:
    $$P(X = x) = \int_x^x f(t) \, dt = 0$$
    Density $f(x)$ is NOT a probability; it is probability per unit interval. Only integrals yield probabilities!
  \end{alertblock}
\end{frame}

% Slide 9: Worked Example 1: The Uniform Distribution
\begin{frame}{Worked Example 1: The Continuous Uniform Distribution}
  Let's define a continuous random variable $X$ distributed uniformly between $a$ and $b$ ($X \sim U(a, b)$).
  \begin{itemize}
    \item \textbf{Analytical PDF:} Since the probability density is constant ($C$) between $a$ and $b$ and zero elsewhere:
      $$f(x) = \begin{cases} C & \text{if } a \le x \le b \\ 0 & \text{otherwise} \end{cases}$$
    \item \textbf{Solving for C (Normalization):}
      $$\int_{-\infty}^{\infty} f(x) \, dx = \int_{a}^{b} C \, dx = C \cdot [x]_a^b = C(b - a) = 1 \implies C = \frac{1}{b - a}$$
    \item \textbf{Probability of a sub-interval $[c, d]$ (where $a \le c < d \le b$):}
      $$P(c \le X \le d) = \int_{c}^{d} \frac{1}{b - a} \, dx = \frac{d - c}{b - a}$$
  \end{itemize}
\end{frame}

% Slide 10: Section 2 Divider
\begin{frame}[plain]
  \vfill
  \centering
  \begin{beamercolorbox}[sep=8pt,center,shadow=true,rounded=true]{palette primary}
    \usebeamerfont{title}\textbf{Section 2: Cumulative Distribution Functions (CDFs)}
  \end{beamercolorbox}
  \vfill
\end{frame}

% Slide 11: The Cumulative Distribution Function (CDF)
\begin{frame}{The Cumulative Distribution Function (CDF)}
  \begin{block}{Definition}
    The Cumulative Distribution Function, $F(x)$, is the probability that the random variable $X$ is less than or equal to a specific value $x$:
    $$F(x) = P(X \le x) = \int_{-\infty}^{x} f(t) \, dt$$
  \end{block}
  \begin{itemize}
    \item \textbf{Visual Intuition:} If the PDF represents the height of a probability curve at $x$, the CDF represents the accumulated area under that curve from the far left up to $x$.
    \item \textbf{Deriving the PDF from the CDF:}
      By the Fundamental Theorem of Calculus, if $F(x)$ is differentiable, we recover the probability density by taking the derivative:
      $$f(x) = \frac{dF(x)}{dx}$$
  \end{itemize}
\end{frame}

% Slide 12: Mathematical Properties of CDFs
\begin{frame}{Mathematical Properties of CDFs}
  Every valid CDF $F(x)$ must satisfy three fundamental properties:
  \begin{enumerate}
    \item \textbf{Monotonicity:} $F(x)$ is non-decreasing. If $x_1 < x_2$, then $F(x_1) \le F(x_2)$. Adding more interval range can never decrease the accumulated probability.
    \item \textbf{Limits at Infinity:}
      $$\lim_{x \to -\infty} F(x) = 0 \quad \text{and} \quad \lim_{x \to \infty} F(x) = 1$$
      As we integrate to the far left, we cover nothing. As we integrate to the far right, we cover the entire sample space.
    \item \textbf{Right-Continuity:} The function is continuous when approached from the right:
      $$\lim_{t \to x^+} F(t) = F(x)$$
  \end{enumerate}
  \begin{exampleblock}{Probability of any interval $[a, b]$ via CDF:}
    $$P(a < X \le b) = F(b) - F(a)$$
    This is extremely useful because we do not have to perform integration every time!
  \end{exampleblock}
\end{frame}

% Slide 13: Worked Example 2: The Exponential Distribution CDF
\begin{frame}{Worked Example 2: The Exponential Distribution CDF}
  The Exponential distribution is widely used to model decay times (e.g., radioactive decay, or the delay times of Type Ia supernovae).
  \begin{itemize}
    \item \textbf{PDF ($x \ge 0$ and decay rate $\lambda > 0$):}
      $$f(x) = \lambda e^{-\lambda x}$$
    \item \textbf{Step-by-Step CDF Derivation:}
      $$F(x) = \int_{-\infty}^{x} f(t) \, dt = \int_{0}^{x} \lambda e^{-\lambda t} \, dt$$
      Recall that the antiderivative of $\lambda e^{-\lambda t}$ is $-e^{-\lambda t}$:
      $$F(x) = \left[ -e^{-\lambda t} \right]_0^x = \left( -e^{-\lambda x} \right) - \left( -e^{-\lambda (0)} \right)$$
      $$F(x) = -e^{-\lambda x} - (-1) = 1 - e^{-\lambda x}$$
    \item \textbf{Final CDF expression:}
      $$F(x) = \begin{cases} 0 & \text{if } x < 0 \\ 1 - e^{-\lambda x} & \text{if } x \ge 0 \end{cases}$$
  \end{itemize}
\end{frame}

% Slide 14: The Problem with PDFs for Observational Data
\begin{frame}{The Problem with PDFs for Observational Data}
  Why do astronomers prefer cumulative functions (CDFs) over density functions (PDFs) when working with real telescope data?
  \begin{itemize}
    \item \textbf{Observational Data is Discrete:} We observe a list of individual galaxy masses $\{M_1, M_2, \dots, M_n\}$, not a continuous equation.
    \item \textbf{The Drawbacks of Histograms (PDF Estimation):}
      \begin{itemize}
        \item \alert{Bin Width Sensitivity:} Bins too narrow $\rightarrow$ noisy and jagged distribution. Bins too wide $\rightarrow$ merges distinct physical features (e.g., washes out the red/blue sequence bimodality).
        \item \alert{Bin Edge Sensitivity:} Shifting the starting edge of the bins (even by a tiny fraction) changes which points fall into which bins, altering the visual shape of the histogram.
        \item \alert{Information Loss:} Binning throws away the exact coordinate value of each object, reducing precision.
      \end{itemize}
    \item \textbf{The Solution:} We need a distribution function that uses the exact coordinates of every data point without grouping them.
  \end{itemize}
\end{frame}

% Slide 15: Section 3 Divider
\begin{frame}[plain]
  \vfill
  \centering
  \begin{beamercolorbox}[sep=8pt,center,shadow=true,rounded=true]{palette primary}
    \usebeamerfont{title}\textbf{Section 3: Empirical Cumulative Distribution Functions (eCDFs)}
  \end{beamercolorbox}
  \vfill
\end{frame}

% Slide 16: The Empirical CDF (eCDF)
\begin{frame}{The Empirical Cumulative Distribution Function (eCDF)}
  For a discrete sample of size $n$, the **Empirical CDF** ($F_n(x)$) is the fraction of data points that are less than or equal to $x$.
  \begin{block}{Mathematical Definition}
    $$F_n(x) = \frac{1}{n} \sum_{i=1}^n \mathbb{I}_{(-\infty, x]}(x_i)$$
    where the indicator function is defined as:
    $$\mathbb{I}_{(-\infty, x]}(x_i) = \begin{cases} 1 & \text{if } x_i \le x \\ 0 & \text{if } x_i > x \end{cases}$$
  \end{block}
  \begin{itemize}
    \item \textbf{Geometry of the eCDF:}
      \begin{itemize}
        \item It is a **step function** (staircase shape).
        \item It starts at exactly $0.0$ on the far left.
        \item At each sorted data value, it jumps upward by exactly $1/n$.
        \item It reaches exactly $1.0$ at the largest data value.
      \end{itemize}
  \end{itemize}
\end{frame}

% Slide 17: Step-by-Step Construction of an eCDF
\begin{frame}{Step-by-Step Construction of an eCDF}
  How do we construct an eCDF from a raw list of data coordinates?
  \begin{enumerate}
    \item \textbf{Collect the Raw Sample:} $\{x_1, x_2, \dots, x_n\}$.
    \item \textbf{Sort in Ascending Order:} Rearrange the sample such that:
      $$x_{(1)} \le x_{(2)} \le x_{(3)} \le \dots \le x_{(n)}$$
      (These are called the *order statistics*).
    \item \textbf{Assign Cumulative Probabilities:}
      At each coordinate $x_{(i)}$, the cumulative probability is:
      $$F_n(x_{(i)}) = \frac{i}{n}$$
    \item \textbf{Write the Piecewise Step Function:}
      $$F_n(x) = \begin{cases} 
        0 & \text{if } x < x_{(1)} \\
        \frac{1}{n} & \text{if } x_{(1)} \le x < x_{(2)} \\
        \vdots & \\
        \frac{i}{n} & \text{if } x_{(i)} \le x < x_{(i+1)} \\
        \vdots & \\
        1 & \text{if } x \ge x_{(n)}
      \end{cases}$$
  \end{enumerate}
\end{frame}

% Slide 18: Worked Example 3: eCDF Calculation by Hand
\begin{frame}{Worked Example 3: eCDF Calculation by Hand}
  Let's build the eCDF for a tiny dataset of galaxy stellar masses ($n = 4$):
  $$S = \{10.5, 9.2, 11.1, 9.8\} \, \log_{10} M_\odot$$
  \begin{enumerate}
    \item \textbf{Sort the data:}
      $$x_{(1)} = 9.2, \quad x_{(2)} = 9.8, \quad x_{(3)} = 10.5, \quad x_{(4)} = 11.1$$
    \item \textbf{Calculate step heights (increments of $1/4 = 0.25$):}
      \begin{itemize}
        \item For $x < 9.2$: $F_4(x) = 0.0$
        \item For $9.2 \le x < 9.8$: $F_4(x) = 0.25$
        \item For $9.8 \le x < 10.5$: $F_4(x) = 0.50$
        \item For $10.5 \le x < 11.1$: $F_4(x) = 0.75$
        \item For $x \ge 11.1$: $F_4(x) = 1.00$
      \end{itemize}
  \end{enumerate}
  \begin{block}{Pedagogical Note}
    Notice that the eCDF preserves every single coordinate measurement perfectly. No binning was used, and we can read the exact median ($10.15 \log_{10} M_\odot$ at $F_4(x)=0.50$) directly from the function.
  \end{block}
\end{frame}

% Slide 19: Theoretical Foundations: Glivenko-Cantelli Theorem
\begin{frame}{Theoretical Foundations: Glivenko-Cantelli Theorem}
  Why are we allowed to use the empirical CDF ($F_n(x)$) as a proxy for the true population CDF ($F(x)$)?
  \begin{block}{The Glivenko-Cantelli Theorem (Fundamental Theorem of Statistics)}
    As the sample size $n$ approaches infinity, the empirical CDF $F_n(x)$ converges uniformly to the true underlying CDF $F(x)$ with probability 1:
    $$P\left( \lim_{n \to \infty} \sup_{x} |F_n(x) - F(x)| = 0 \right) = 1$$
  \end{block}
  \begin{itemize}
    \item \textbf{What this means physically:}
      If we observe enough galaxies, our step-like staircase eCDF will perfectly smooth out and become identical to the true physical distribution function of the universe.
    \item This theorem justifies using the difference between two eCDFs to test if they come from the same parent population.
  \end{itemize}
\end{frame}

% Slide 20: Section 4 Divider
\begin{frame}[plain]
  \vfill
  \centering
  \begin{beamercolorbox}[sep=8pt,center,shadow=true,rounded=true]{palette primary}
    \usebeamerfont{title}\textbf{Section 4: The Hypothesis Testing Paradigm}
  \end{beamercolorbox}
  \vfill
\end{frame}

% Slide 21: The Logic of Hypothesis Testing
\begin{frame}{The Logic of Hypothesis Testing}
  Hypothesis testing is a rigorous framework for making decisions under uncertainty:
  \begin{itemize}
    \item \textbf{Null Hypothesis ($H_0$):} The ``status quo'' or ``no effect'' statement. We assume it is true by default.
      \begin{itemize}
        \item \textit{Example:} ``Supernova host galaxies have the exact same mass distribution as normal background galaxies.''
      \end{itemize}
    \item \textbf{Alternative Hypothesis ($H_1$):} The scientific claim we wish to establish.
      \begin{itemize}
        \item \textit{Example:} ``Supernova host galaxies have systematically different masses due to progenitor physics.''
      \end{itemize}
  \end{itemize}
  \begin{block}{The Courtroom Analogy}
    The null hypothesis is like the presumption of innocence:
    $$\text{Assume } H_0 \text{ is true} \rightarrow \text{Gather Data} \rightarrow \text{Is the data highly improbable under } H_0? \rightarrow \text{Reject } H_0$$
    If the data is NOT highly improbable, we do NOT ``prove $H_0$ is true''; we simply state we ``fail to reject $H_0$'' (verdict: Not Guilty, which is not the same as Innocent).
  \end{block}
\end{frame}

% Slide 22: Test Statistics, Critical Regions, & Significance
\begin{frame}{Test Statistics, Critical Regions, \& Significance}
  How do we mathematically decide when to reject the Null Hypothesis?
  \begin{itemize}
    \item \textbf{Test Statistic ($T$):} A single mathematical value computed from our sample data that measures the deviation from the null hypothesis.
    \item \textbf{Significance Level ($\alpha$):} The threshold probability we set *before* the experiment for rejecting $H_0$. By scientific convention, we set:
      $$\alpha = 0.05 \quad (5\% \text{ threshold})$$
    \item \textbf{Critical Region:} The set of all test statistic values that are highly unlikely to occur if $H_0$ is true.
    \item \textbf{Critical Value ($T_{\text{crit}}$):} The boundary of the critical region. If our calculated test statistic $T \ge T_{\text{crit}}$, we reject $H_0$.
  \end{itemize}
\end{frame}

% Slide 23: Type I and Type II Errors
\begin{frame}{Type I and Type II Errors}
  Because we work with random samples, there is always a chance we make the wrong decision:
  \begin{table}
    \centering
    \begin{tabular}{lll}
      \toprule
      Decision & $H_0$ is Actually True & $H_0$ is Actually False \\
      \midrule
      \textbf{Fail to Reject $H_0$} & Correct Decision ($1-\alpha$) & Type II Error ($\beta$) \\
      \textbf{Reject $H_0$ (Significant)} & \alert{Type I Error ($\alpha$)} & Correct Decision (Power: $1-\beta$) \\
      \bottomrule
    \end{tabular}
  \end{table}
  \begin{itemize}
    \item \textbf{Type I Error ($\alpha$):} False Positive. We claim supernova hosts are different when they are actually identical. We control this directly by setting $\alpha = 0.05$.
    \item \textbf{Type II Error ($\beta$):} False Negative. We fail to detect a real physical difference because our sample is too small or noisy.
    \item \textbf{Statistical Power ($1-\beta$):} The probability of correctly rejecting a false null hypothesis. To increase power, we must increase our sample size $n$.
  \end{itemize}
\end{frame}

% Slide 24: Demystifying the p-value
\begin{frame}{Demystifying the p-value}
  \begin{block}{What is a p-value?}
    The probability, assuming the null hypothesis is true, of obtaining a test statistic at least as extreme as the one calculated from our observed sample:
    $$\text{p-value} = P(T \ge T_{\text{obs}} \mid H_0)$$
  \end{block}
  \begin{itemize}
    \item \alert{CRITICAL WARNING:} The p-value is NOT the probability that the null hypothesis is true, i.e., $P(H_0 \mid \text{Data}) \ne \text{p-value}$.
    \item It is a conditional probability in the opposite direction: $P(\text{Data} \mid H_0)$.
    \item \textbf{Decision Rule:}
      \begin{itemize}
        \item If $\text{p-value} < \alpha$: The result is highly unlikely under the null. \textbf{Reject $H_0$}.
        \item If $\text{p-value} \ge \alpha$: The result is consistent with random chance. \textbf{Fail to Reject $H_0$}.
      \end{itemize}
  \end{itemize}
\end{frame}

% Slide 25: Section 5 Divider
\begin{frame}[plain]
  \vfill
  \centering
  \begin{beamercolorbox}[sep=8pt,center,shadow=true,rounded=true]{palette primary}
    \usebeamerfont{title}\textbf{Section 5: One-Dimensional Distribution Tests}
  \end{beamercolorbox}
  \vfill
\end{frame}

% Slide 26: The 1D Kolmogorov-Smirnov (K-S) Test
\begin{frame}{The 1D Kolmogorov-Smirnov (K-S) Test}
  The K-S test evaluates whether two independent continuous samples are drawn from the same underlying distribution.
  \begin{itemize}
    \item Let $F_1(x)$ be the eCDF of Sample 1 (size $n_1$) and $F_2(x)$ be the eCDF of Sample 2 (size $n_2$).
    \item \textbf{The K-S Statistic ($D$):} The maximum vertical distance between the two step functions:
      $$D = \sup_x |F_1(x) - F_2(x)|$$
      Where $\sup$ (supremum) represents the absolute maximum separation.
    \item \textbf{Double-Bound Evaluation (Technical Detail):}
      Because eCDFs are step functions, the maximum distance must be checked at both the left-hand limit and right-hand limit of every data point:
      $$D = \max_{i} \left( \max \left( |F_1(x_i) - F_2(x_i)|, |F_1(x_i) - F_2(x_{i-1})| \right) \right)$$
    \item Bounded strictly between $0.0$ (identical eCDFs) and $1.0$ (complete separation).
  \end{itemize}
\end{frame}

% Slide 27: Worked Example 4: Manual 1D K-S Computation
\begin{frame}{Worked Example 4: Manual 1D K-S Computation}
  Let's calculate $D$ for two small samples ($n_1=3, n_2=3$):
  \begin{itemize}
    \item \textbf{Sample 1 ($S_1$):} $\{1.0, 3.0, 5.0\}$
    \item \textbf{Sample 2 ($S_2$):} $\{2.0, 4.0, 6.0\}$
  \end{itemize}
  Let's evaluate the eCDFs $F_1(x)$ and $F_2(x)$ (jumps of $1/3 \approx 0.33$):
  \begin{table}
    \centering
    \begin{tabular}{cccc}
      \toprule
      $x$ & $F_1(x)$ (Sample 1) & $F_2(x)$ (Sample 2) & Absolute Difference $|F_1(x) - F_2(x)|$ \\
      \midrule
      1.0 & 0.33 & 0.00 & 0.33 \\
      2.0 & 0.33 & 0.33 & 0.00 \\
      3.0 & 0.67 & 0.33 & \textbf{0.34} \\
      4.0 & 0.67 & 0.67 & 0.00 \\
      5.0 & 1.00 & 0.67 & 0.33 \\
      6.0 & 1.00 & 1.00 & 0.00 \\
      \bottomrule
    \end{tabular}
  \end{table}
  The maximum vertical separation is $D = 0.34$ (occurring at $x=3.0$).
\end{frame}

% Slide 28: The Fatal Flaw of K-S: Tail Blindness
\begin{frame}{The Fatal Flaw of K-S: Tail Blindness}
  While the K-S test is powerful, it has a major mathematical limitation:
  \begin{itemize}
    \item All cumulative distributions must start at $0.0$ on the left and end at $1.0$ on the right:
      $$\lim_{x \to -\infty} |F_1(x) - F_2(x)| = 0 \quad \text{and} \quad \lim_{x \to \infty} |F_1(x) - F_2(x)| = 0$$
    \item Therefore, the variance of the difference between the two eCDFs shrinks to zero at both extremes.
    \item The K-S statistic $D$ is almost always found near the **median** ($F(x) \approx 0.5$) of the distributions.
    \item \alert{The Physical Consequence:}
      If supernova host galaxies have the same average median mass as field galaxies, but differ at the tails (e.g., they have more dwarf host galaxies or more massive host galaxies), the K-S test will be blind to this difference and fail to reject $H_0$.
  \end{itemize}
\end{frame}

% Slide 29: The Anderson-Darling (A-D) Test
\begin{frame}{The Anderson-Darling (A-D) Test}
  The A-D test solves the tail-blindness problem by introducing a weighting function that divides the squared difference by the variance of the pooled distribution.
  \begin{block}{Mathematical Formulation}
    $$A^2 = n_1 n_2 \int_{-\infty}^{\infty} \frac{[F_1(x) - F_2(x)]^2}{H(x)(1 - H(x))} \, dH(x)$$
    Where:
    \begin{itemize}
      \item $H(x)$ is the **pooled cumulative distribution** of both samples:
        $$H(x) = \frac{n_1 F_1(x) + n_2 F_2(x)}{n_1 + n_2}$$
      \item The denominator term $H(x)(1 - H(x))$ represents the variance of this pooled distribution.
    \end{itemize}
  \end{block}
\end{frame}

% Slide 30: Mechanics of the A-D Weight Function
\begin{frame}{Mechanics of the A-D Weight Function}
  Let's look at the weight term $W(H) = \frac{1}{H(x)(1-H(x))}$:
  \begin{itemize}
    \item At the center of the distribution ($H(x) \approx 0.5$):
      $$W(0.5) = \frac{1}{0.5 \cdot 0.5} = 4.0$$
    \item At the extreme tails of the distribution ($H(x) \approx 0.01$ or $0.99$):
      $$W(0.01) = \frac{1}{0.01 \cdot 0.99} \approx 101.0$$
    \item The squared differences at the tails are multiplied by a weighting factor that is over **25 times larger** than at the center!
    \item By integrating over the entire range with this tail-weighting function, the A-D test becomes highly sensitive to differences at the extremes.
  \end{itemize}
\end{frame}

% Slide 31: Discrete Summation Formula for A-D
\begin{frame}{Discrete Summation Formula for A-D}
  For computational implementation, SciPy converts the continuous integral into a summation over the combined, sorted sample of size $N = n_1 + n_2$:
  \begin{block}{Two-Sample A-D Summation}
    $$A^2 = \frac{n_1 n_2}{N} \sum_{i=1}^{N-1} \frac{\left( F_1(z_i) - F_2(z_i) \right)^2}{H(z_i)(1 - H(z_i))}$$
    Where:
    \begin{itemize}
      \item $z_i$ are the elements of the combined and sorted dataset.
      \item $F_1(z_i)$ and $F_2(z_i)$ are the eCDFs of the individual samples evaluated at $z_i$.
      \item $H(z_i) = i / N$ is the pooled eCDF step height.
    \end{itemize}
  \end{block}
  Because the weights blow up to infinity if $H(z_i) = 0$ or $1$, the sum is evaluated from $i=1$ to $N-1$, excluding the boundaries.
\end{frame}

% Slide 32a: Side-by-Side Comparison: K-S vs. A-D Matrix
\begin{frame}{Side-by-Side Comparison: K-S vs. A-D Matrix}
  \begin{table}
    \centering
    \begin{tabular}{lll}
      \toprule
      Feature & Kolmogorov-Smirnov (K-S) & Anderson-Darling (A-D) \\
      \midrule
      \textbf{Sensitivity} & Center of distribution (median) & Tails of distribution (extremes) \\
      \textbf{Mathematical Metric} & Supremum vertical distance ($D$) & Integrated weighted variance ($A^2$) \\
      \textbf{Critical Values} & Analytical formula & Numerical simulations \\
      \textbf{Stability} & Highly stable, robust to noise & Sensitive to individual extreme noise \\
      \textbf{Astrophysical Use} & Finding overall median offsets & Catching dwarf or massive outliers \\
      \bottomrule
    \end{tabular}
  \end{table}
\end{frame}

% Slide 32b: Side-by-Side Comparison: Methodological Guidelines
\begin{frame}{Side-by-Side Comparison: Methodological Guidelines}
  \begin{alertblock}{Methodological Rule of Thumb}
    Always run both tests!
    \begin{itemize}
      \item If K-S is significant but A-D is not: The populations differ in their bulk properties.
      \item If A-D is significant but K-S is not: The populations differ only in their extreme outliers.
    \end{itemize}
    \alert{Critical Note:} SciPy's \texttt{anderson\_ksamp} returns a normalized p-value capped at $0.001 \le p \le 0.25$. Do not panic if your p-value is exactly $0.001$ or $0.25$; this is a limitation of the interpolation tables!
  \end{alertblock}
\end{frame}

% Slide 33: Section 6 Divider
\begin{frame}[plain]
  \vfill
  \centering
  \begin{beamercolorbox}[sep=8pt,center,shadow=true,rounded=true]{palette primary}
    \usebeamerfont{title}\textbf{Section 6: Multi-Dimensional Statistics \& Joint Distributions}
  \end{beamercolorbox}
  \vfill
\end{frame}

% Slide 34: Galaxy Evolution in 2D Space
\begin{frame}{Galaxy Evolution in 2D Space}
  Galaxies do not evolve in one-dimensional isolation. Their physical parameters are tightly coupled.
  \begin{itemize}
    \item \textbf{The Stellar Mass--SFR Correlation:}
      \begin{itemize}
        \item Active, star-forming galaxies lie along a diagonal corridor called the \textbf{Star-Forming Main Sequence} (SFMS).
        \item Quenched, dead elliptical galaxies lie far below it in the \textbf{Red Sequence}.
      \end{itemize}
    \item \alert{The Danger of 1D Testing:}
      If we run a 1D test on Mass, we lose all star-formation information. If we test star formation, we ignore mass.
    \item A population could shift diagonally (along the correlation line) or rotate in the plane. 1D tests on individual axes would fail to detect this change!
    \item We need a joint probability distribution:
      $$F(x, y) = P(X \le x, Y \le y) = \int_{-\infty}^{x} \int_{-\infty}^{y} f(u, v) \, dv \, du$$
  \end{itemize}
\end{frame}

% Slide 35: The 2D Cumulative Ordering Problem
\begin{frame}{The 2D Cumulative Ordering Problem}
  Why is generalizing the K-S test to 2D mathematically challenging?
  \begin{itemize}
    \item In 1D, numbers have a natural ordering: we can only integrate from left to right ($X \le x$).
    \item In 2D, there is no unique ordering. An evaluation point $(x_0, y_0)$ splits the plane into **four quadrants**, and we can integrate in four possible directions:
      \begin{align*}
        Q_1(x_0, y_0) = P(X \le x_0, Y \le y_0) & \quad (\text{Bottom-Left}) \\
        Q_2(x_0, y_0) = P(X > x_0, Y \le y_0) & \quad (\text{Bottom-Right}) \\
        Q_3(x_0, y_0) = P(X \le x_0, Y > y_0) & \quad (\text{Top-Left}) \\
        Q_4(x_0, y_0) = P(X > x_0, Y > y_0) & \quad (\text{Top-Right})
      \end{align*}
    \item A simple 2D CDF that only integrates in one direction (e.g., $Q_1$) ignores the data structure in the other three directions.
  \end{itemize}
\end{frame}

% Slide 36: The 2D K-S Test (Fasano \& Franceschini 1987)
\begin{frame}{The 2D K-S Test (Fasano \& Franceschini 1987)}
  Fasano \& Franceschini resolved the ordering problem by evaluating the differences in all four quadrants around every data point.
  \begin{itemize}
    \item Let the plane be split into four quadrants around any coordinate $(x_i, y_i)$ in our combined sample.
    \item For each quadrant $q \in \{1, 2, 3, 4\}$, we compute the fraction of points from Sample 1 ($q^{(1)}$) and Sample 2 ($q^{(2)}$) that fall inside it.
    \item \textbf{The 2D K-S Statistic ($D_{2\text{D}}$):} The maximum absolute difference across all four quadrants, evaluated at every data point in the combined sample:
      $$D_{2\text{D}} = \max_{(x_i, y_i)} \left( \max_{q \in \{1,2,3,4\}} |q^{(1)} - q^{(2)}| \right)$$
    \item \textbf{Computational Complexity:} Evaluating all quadrants around all $N$ data points requires $O(N^2)$ calculations.
  \end{itemize}
\end{frame}

% Slide 37: Peacock's Significance & Correlation Correction
\begin{frame}{Peacock's Significance \& Correlation Correction}
  Because coordinates in 2D space are often physically correlated (like Mass and SFR), the standard 1D significance tables cannot be used. We must apply Peacock's numerical approximation:
  \begin{enumerate}
    \item \textbf{Calculate the Correlation Coefficient ($r$):} Compute the average Pearson correlation coefficient of both samples:
      $$r = \frac{1}{2}(r_1 + r_2)$$
    \item \textbf{Calculate the Normalized Test Statistic ($\lambda$):}
      $$\lambda = \frac{D_{2\text{D}} \sqrt{N_{\text{eff}}}}{1 - 0.5 r^2}, \quad \text{where } N_{\text{eff}} = \frac{n_1 n_2}{n_1 + n_2}$$
      The term $1 - 0.5r^2$ acts as a correction factor: if the variables are highly correlated ($r \to 1$), the denominator shrinks, amplifying $\lambda$ to account for the reduced degrees of freedom.
    \item \textbf{Calculate the p-value:}
      $$P \approx 2 \cdot \exp\left(-2(\lambda - 0.5)^2\right)$$
      Clamped strictly to the interval $[0.0, 1.0]$.
  \end{enumerate}
\end{frame}

% Slide 38: Interpreting Our Task 9 Results
\begin{frame}{Interpreting Our Task 9 Results}
  When you execute the code in the solution notebook, you compare supernova host galaxies with background galaxies:
  \begin{itemize}
    \item \textbf{Core-Collapse SNe:} Host galaxies cluster tightly along the active Star-Forming Main Sequence.
      \begin{itemize}
        \item CC-SNe trace massive star formation, so they are absent in quenched ellipticals.
      \end{itemize}
    \item \textbf{Type Ia SNe:} Host galaxies are bimodal, split between star-forming spirals and quenched ellipticals.
      \begin{itemize}
        \item SNe Ia are driven by white dwarfs, which have long delay times.
      \end{itemize}
    \item \textbf{Statistical Power Comparison:}
      \begin{itemize}
        \item 1D K-S on Stellar Mass: $p \approx 10^{-3}$
        \item 1D K-S on sSFR: $p \approx 10^{-4}$
        \item \alert{2D K-S on Mass \& sSFR:} $p \approx 10^{-6}$
      \end{itemize}
    \item The 2D test is orders of magnitude more significant because it captures the correlation structure (SFMS vs. Red Sequence) that 1D tests wash out!
  \end{itemize}
\end{frame}

% Slide 39: Summary \& Homework
\begin{frame}{Summary \& Meet 2 Preview}
  \begin{block}{Key Takeaways}
    \begin{itemize}
      \item CDFs integrate PDFs, smoothing out statistical noise and preserving exact measurements.
      \item K-S statistics evaluate the maximum vertical distance between eCDFs.
      \item A-D statistics weight the tails, catching differences at distribution extremes.
      \item The 2D K-S test evaluates four quadrants to capture physical correlations in multi-dimensional space.
    \end{itemize}
  \end{block}
  \begin{description}
    \item[Homework:] Open \texttt{stellar\_trace\_assignment2.ipynb} and complete the missing code sections for Tasks 7--9. Verify your code compiles and reproduces the 1D/2D statistics.
    \item[Meet 2 Preview:] Next meet, we will address class imbalances using SMOTE under strict train-test separation to prevent data leakage, and train our first Multi-Layer Perceptron classifier.
  \end{description}
\end{frame}

\end{document}
